\section{System Architecture}
\label{sec:architecture}

Durable Swarms coordinates independent, concurrent workflows over a fixed pool
of virtual machines (VMs). Our objective is to preserve completed work across
worker failures while separating application logic from placement, monitoring,
and recovery. We instantiate the application layer with simulated workflows
that model a leader and its subtasks. This choice allows us to examine execution
and recovery mechanisms without depending on live language-model agents; it
does not model semantic agent misalignment or establish the quality of generated
code.

\subsection{Execution Model and Components}
\label{sec:components}

A workload consists of one logical leader and a set of subtask workers. Each
worker executes an ordered sequence of steps, and dependencies between subtasks
form a directed acyclic graph. The current specification declares the subtasks,
their dependencies, and expected results at submission time. Steps represent
compute delays or simulated side effects. A seed determines any sampled delays;
the normalized specification and fault schedule are persisted with the workload.
These inputs are reproducible, although concurrent scheduling and elapsed
execution time remain timing-dependent.

Figure~\ref{fig:architecture} shows the three runtime components: a single
orchestrator, a durable store, and VM daemons that host disposable worker
processes. The application submits a specification and retrieves status and
results through an HTTP interface. It does not assign workers to machines or
reconstruct them after a failure. The prototype exposes a declarative workflow
interface; a general interface for programmer-supplied execution functions
remains future work.

\begin{figure}[t]
  \centering
  \begin{tikzpicture}[
    box/.style={draw, rounded corners=2pt, align=center,
                minimum height=0.85cm, font=\small},
    link/.style={<->, >=Stealth, semithick},
    every node/.style={font=\small}
  ]
    \node[box, text width=2.35cm] (app) at (0,2.7) {Application};
    \node[box, text width=2.35cm] (orch) at (0,1.1) {Orchestrator};
    \node[box, text width=2.0cm] (store) at (4.15,1.1)
      {Durable store\\Log and receipts};
    \node[box, text width=3.2cm] (pool) at (0,-0.8)
      {Fixed VM pool\\Daemons and worker slots};
    \draw[link] (app) -- node[right, align=left]
      {Tasks, status,\\results} (orch);
    \draw[link] (orch) -- node[above] {State} (store);
    \draw[link] (orch) -- node[right, align=left]
      {Assignments\\Worker reports} (pool);
  \end{tikzpicture}
  \caption{Durable Swarms separates disposable worker execution from durable
  coordination state. Worker updates and simulated effects pass through the
  orchestrator, which validates the current attempt before committing changes.}
  \label{fig:architecture}
\end{figure}

Each VM has a configured number of worker slots. Its daemon registers a session,
polls for assignments, launches worker subprocesses, and reports their liveness.
The orchestrator owns placement and infrastructure failure detection. A leader
creates its declared children through the orchestrator, collects their results,
and authorizes retries for failed children. Thus, infrastructure detection is
centralized while child-recovery decisions remain associated with the parent
workload. Recovery of the leader itself is initiated by the orchestrator.

\subsection{Placement and Coordination}
\label{sec:placement}

The scheduler assigns workers whose dependencies have completed to available
slots on live VMs. Within a scheduling cycle, it considers at most one ready
worker per workload per pass and repeats while capacity remains. Placement
prefers a VM other than the worker's previous failed placement, then a VM with
lower slot occupancy relative to its capacity. Migration is therefore preferred
when another suitable VM is available, rather than required for every recovery.
Slot accounting bounds assigned concurrency; it does not measure or enforce
per-worker CPU or memory consumption.

A leader that has completed its own steps releases its slot while awaiting
child results. Child failure or completion of all children makes the leader
runnable again. This prevents a waiting leader from occupying the only slot
needed to execute its subtasks. Completion requires every declared worker to
finish and the leader to aggregate their results.

\subsection{Durable State}
\label{sec:durable-state}

The durable store uses SQLite with write-ahead logging and full synchronization.
An event records its sequence number, workload identifier, timestamp, state
version, event type, and payload. Worker-specific events additionally identify
the logical worker, parent, and execution attempt. The history includes
assignments, progress, subtask creation, effect intents and receipts, completion,
and recovery transitions. State-changing workload transactions also append a
snapshot containing worker state and each worker's next incomplete step.

The workload table is a materialized projection of these snapshots. On
orchestrator restart, the projection is reconstructed from the checkpoint
history. Changes to a workload, its events, and its snapshot commit in one
transaction before acknowledgment. Worker VMs therefore need no durable local
execution state. Recovery skips completed steps, but an interrupted compute
delay may execute again because its completion has not yet been checkpointed.

\subsection{Failure Detection and Recovery}
\label{sec:recovery}

We distinguish VM liveness from worker progress. A VM heartbeat renews the VM
lease and only those unexpired worker attempts that its daemon reports as
alive. A separate progress timeout detects workers that remain alive but stop
advancing. Repeated reports of the same progress value do not reset this timer.
Timeouts can also suspect delayed or disconnected workers, so recovery must
remain safe even when an old process continues running.

Each assignment receives a monotonically increasing attempt identifier.
Worker requests carry the workload and worker identifiers, attempt identifier,
VM identifier, and VM session. Before accepting a mutation, the orchestrator
checks this identity against the current assignment and verifies that both
leases remain valid. Superseded attempts cannot update progress, commit
checkpoints, or apply simulated effects through the service.

When a failure is detected, the orchestrator invalidates the current attempt
and retains its durable step cursor. A failed leader is queued for replacement;
a failed child awaits retry authorization from its leader. The scheduler then
assigns a replacement attempt using the retained checkpoint. Invalid final
results fail the workload, while exhaustion of the configured failure budget
escalates it. These outcomes make unsuccessful recovery explicit rather than
reporting an incomplete workload as successful.

\subsection{Replay and Side-Effect Safety}
\label{sec:effects}

Checkpointing alone cannot prevent duplicate effects: a worker may apply an
operation and fail before recording step completion. We therefore give each
simulated effect a stable identifier,
$o = H(w, a, s)$, derived from workload identifier $w$, logical worker identifier
$a$, and step index $s$. The identifier excludes the execution attempt and VM,
so a replacement worker refers to the same logical operation.

Execution proceeds in three stages. First, the worker records a durable intent
for the effect. Second, the service applies the simulated mutation and stores
its receipt atomically; a repeated operation identifier returns the original
receipt. Third, the worker advances its step checkpoint after the receipt is
durable. A failure before mutation leaves the operation unapplied. A failure
after mutation but before checkpointing causes the replacement to retrieve the
existing receipt without repeating the mutation. Once the checkpoint commits,
the completed step is skipped during recovery.

Simulated effects increment counters scoped to each workload. The
counter mutation, receipt, and completion event share a SQLite transaction,
which closes the ambiguity window for this simulated destination. This provides
at-least-once step execution with idempotent effect application, not a general
exactly-once guarantee for arbitrary external tools. An external destination
would require its own idempotency or reconciliation mechanism. Inconsistent
effect records detected by the service cause escalation rather than automatic
reapplication.

\subsection{Scope and Current Tradeoffs}
\label{sec:architecture-scope}

The design assumes a trusted private network, one active orchestrator, and
durable storage that survives process restarts. Worker crashes, VM-daemon loss,
stalls, retry loops, invalid results, and interruption after an effect are
represented by the simulator. Loss of the orchestrator's underlying storage and
Byzantine workers are outside the recovery model. Resetting a VM registration
fences its assignments and permits a new daemon session; provisioning or booting
VMs remains the responsibility of the existing VM manager.

The implementation favors explicit recovery state over optimized metadata
handling. Progress and lease updates can write complete workload snapshots,
and scheduling and monitoring scan workload records. Scheduling passes also
restart in insertion order, so the current policy does not establish fairness
across successive cycles. Reducing snapshot frequency, indexing active work,
and retaining a rotating scheduling cursor are targeted refinements for the
resource-efficiency goal. We do not claim large-cluster scalability or measured
efficiency from the present architecture alone.
